\section{性爱前后的清洁、准备与经期性交}

\subsection{事前准备}

\begin{itemize}
	\item \textbf{洗手、修剪指甲}：指甲过长或边缘粗糙容易划伤阴道黏膜与外阴皮肤，事前修剪、磨平即可避免；
	\item \textbf{清洗外阴即可，不必冲洗阴道}：用清水清洗外阴、阴茎与包皮内侧即可。阴道具有\textbf{自净功能}，用洗液冲洗阴道内部反而会破坏正常菌群、增加感染风险；
	\item \textbf{提前如厕}：排空膀胱与肠道，可减少过程中的尿意、胀气与意外；
	\item \textbf{备好物品}：避孕套、润滑剂、纸巾或毛巾放在随手可及处，避免临时中断扫兴。
\end{itemize}

\subsection{事后清洁}

\begin{itemize}
	\item \textbf{双方清洗外生殖器}：清水冲洗外阴、阴茎即可，无需使用刺激性洗液；
	\item \textbf{尽快排尿}：\textbf{女性性交后应尽快排一次尿}，可显著降低尿路感染（俗称"蜜月膀胱炎"）的风险；男性排尿也有助于冲刷尿道、减少感染；
	\item \textbf{检查安全套}：取出后检查有无破损，打结后丢弃，不要重复使用；
	\item \textbf{不要冲洗阴道内部}，也不要依赖各种"私处护理液"或香氛产品——它们扰乱了原本的酸碱平衡。
\end{itemize}

\subsection{经期可以性交吗？}

\noindent\textbf{医学上的答案是可以}。月经期性交并非禁忌，没有证据表明它本身有害，前提是\textbf{双方自愿、注意卫生}。

\begin{itemize}
	\item \textbf{可能的益处}：部分女性经期盆腔充血、性欲反而升高；对痛经者而言，高潮时释放的内啡肽可能带来一定缓解；
	\item \textbf{需注意的风险}：经期宫颈口略开、子宫内膜正在脱落，感染风险相对略高，因此\textbf{建议使用安全套}——既降低性传播感染风险，也减少血渍、便于清理；
	\item \textbf{避孕误区}：经期\textbf{不是"安全期"}。周期较短或经期较长者，仍可能在经期末尾接近排卵，无保护性交仍有受孕可能，切勿以"经期不会怀孕"为由省去避孕；
	\item \textbf{尊重意愿}：若一方不愿或感到不适，应坦然沟通、彼此体谅。可用深色床单、毛巾或选择在浴室进行，减少顾虑。
\end{itemize}

\begin{tcolorbox}[colback=pink!5!white,colframe=red!50!black,title=贴心小叮咛]
"事前清洁、事后排尿"看似小事，却是预防泌尿生殖道感染的简单有效之举。养成良好的性卫生习惯，是对自己和伴侣身体的体贴。
\end{tcolorbox}
